Consider an employer seeking to hire new employees. 
Clearly, the employer would like to hire the best employees for the task, 
but how will she know which are best fit? 
Typically, the employee will gather information on each candidate,
including their education, work history, reference letters, 
and for many jobs, they will actively conduct interviews. 
Altogether, this information can be viewed as the \emph{signal} available to the employer.


Suppose that candidates can be either \emph{skilled} or \emph{unskilled}. 
If the firm hires an ``unskilled'' candidate, 
it will incur a significant cost on account of lost productivity.
For this reason, the employer would like 
to minimize the number of \emph{False Positive} mistakes, 
instances where \emph{unskilled} candidates are hired. 
% This is one side of the tradeoff. 
On the other hand, the employer desires 
not to \emph{overspend} on the hiring process, 
% e.g., by 
limiting the number of interviews per hired candidate 
(either on average, or absolutely).
However, fewer interviews weakens the signal, 
causing the employer to make more mistakes.
%
At the heart of our model is this inherent trade-off 
between the quality of the signal 
and the cost of obtaining the signal.
This marks a departure from the classical economics literature,
in which the signal is commonly regarded as a given.


Complicating matters, hiring efficiency is not the only desiderata at play. 
In society, employees belong to various \emph{demographic groups}, 
and we may strive to design policies 
that are in some sense \emph{fair} vis-a-vis group membership. 
While \emph{fairness} can be an elusive notion, 
regulators must translate it to concrete rules and laws.
% For example, 
In the United States, a body of anti-discrimination law
dating to the Civil Rights act of 1964, 
subjects decisions that result in disparate outcomes 
(as delineated by race, age, gender, religion, etc)
to extra scrutiny: employers must not only show 
that preference for some groups over others 
did not drive the decision (disparate treatment doctrine) 
but also justify that any observed disparities 
arise from a business necessity (disparate impact doctrine), 
whether or not those disparities were intentional.


In this paper, we seek to understand how a complex hiring process 
would interact with the requirements of fairness.  
We extend the theory on candidate screening and statistical discrimination,
addressing the setting in which employers can subject 
employees to multiple tests,
which we assume to be conditionally independent
given the worker's skill level and group membership.
To build intuition, most of our analysis focuses 
on a Bernoulli model of both worker skill and screening.
Additionally, we begin to extend the traditionally-studied 
Gaussian skill and screening models to the multi-test setting (Section \ref{sec:gaussian}). 


Unlike the classical papers, 
in which an employer's hiring policy is given 
by a simple thresholding rule,
our setting requires greater care 
to derive the optimal employer policy. 
In our setting, we imagine that the employer 
wishes to minimize the number of tests performed
subject to a constraint upper-bounding the false positive rate.
We characterize the optimal policy in this case as a randomized threshold policy.

We also consider the setting in which employers can allocate tests dynamically, 
deciding after each result whether to 
(i) hire the candidate;
(ii) reject the candidate and move on to the next one; or 
(iii) administer a subsequent test. 
In the Bernoulli case, the optimal policy consists 
of administering tests until each candidate's posterior likelihood 
of being a high-skilled worker either 
dips below the prior or rises above a threshold 
determined by the tolerable false positive rate.
We demonstrate that the analysis of this process
can be reduced to a random walk over the log posterior odds
and derive the solution via the corresponding Gambler's ruin problem.

Finally, 
we consider the ramifications for fairness within our model 
when employees, despite possessing similarly-distributed skills, 
are evaluated with differing noise levels.

\subsection{Related work}

The classical economics literature 
on discrimination in employment 
can broadly be divided into two focuses.
The \emph{taste-based discrimination} model 
due to \cite{becker1957economics}
models the market outcomes in a setting 
where employers express an explicit preference 
for hiring members of one group,
acting as if an employee's demographic membership provides utility.
This preference for certain groups
induces a sorting of employees from the disadvantaged group
towards those employers who discriminate the least
with wages ultimately determined by the \emph{marginal discriminator}.
%
Subsequently, \cite{phelps1972statistical} suggested
a statistical mechanism by which similarly-skilled employees 
from different groups might experience differential outcomes: 
the comparative difficulty of screening from one group vs. another.
Many subsequent works extend this analysis,
typically focusing on Gaussian models of worker quality
and conditionally-Gaussian test scores 
\cite{arrow1973theory, aigner1977statistical}.
These papers consider the setting 
where workers are assessed via a single test
characterized by a group-dependent noise level.
Our work is differentiated from these 
by considering richer mechanisms for acquiring signal.

In the more recent literature on fairness in machine learning,
researchers often focus on binary classification, 
with employees characterized by a protected characteristic (group membership),
and other (non-protected) covariates \citep{pedreshi2008discrimination, kamiran2010discrimination, kamishima2011fairness}.
There, the predictor is presumably used 
to guide a consequential decision,
such as allocating some economic good (loans, jobs, etc.) \citep{corbett2018measure}
or assessing some penalty (e.g. risk scores to guide bail decisions) \citep{chouldechova2017fair}.
Papers then focus on various interventions 
for ensuring accurate prediction subject to various constraints 
such as demographic parity (outcomes independent of group membership), 
blindness (model cannot observe group membership), 
and equalized false negative and/or false positive rates \cite{hardt2016equality}. 
Several simple impossibility results 
preclude simultaneously satisfying 
several combinations of these parities 
\cite{berk2017fairness, chouldechova2017fair, kleinberg2016inherent}.
More recently, a number of papers 
have drawn inspiration from economic modeling, 
extending the literature on fairness in classification 
to consider longer-term dynamics, equilibria, 
and the emergence of feedback loops 
\cite{hu2018short, hardt2016equality, ensign2017runaway}. 
Finally, \cite{Verma2018} provide a survey of definitions 
from the algorithmic fairness literature.
